\chapter{1998年上海卷（文）}
\section{填空题}

\begin{problem}
\(\lg20 + \log_{100}25 = \)\fillinblank{}.
\end{problem}

\begin{answer}
2
\end{answer}

\begin{problem}
若函数 \(y = 2\sin x + \sqrt{a} \cos x + 4 \) 的最小值为 \(1\)，则 \(a=\)\fillinblank{}.
\end{problem}

\begin{answer}
5
\end{answer}

\begin{problem}
若 \(\displaystyle\lim_{x\to1}\frac{x^{2}+ax+3}{x^{3}+3}=2 \)，则 \(a=\)\fillinblank{}.
\end{problem}

\begin{answer}
4
\end{answer}

\begin{problem}
函数 \(f(x)=(x-1)^{\frac{1}{3}}+2 \) 的反函数是 \(f^{-1}(x)= \)\fillinblank{}.
\end{problem}

\begin{answer}
\((x-2)^3+1\)，\(x\in\R\)
\end{answer}

\begin{problem}
棱长为\(2\)的正四面体的体积为\fillinblank{}.
\end{problem}

\begin{answer}
\(\dfrac{2\sqrt{2}}{3}\)
\end{answer}

\begin{problem}
某工程的工序流程图如下（工时数单位：天）：

\bitmapfigure[width=0.72\linewidth]{img/1998/shanghai_liberal/q6_fig1.png}

则工程总时数为\fillinblank{}\(\text{天}\).
\end{problem}

\begin{answer}
11
\end{answer}

\begin{problem}
袋内装有大小相同的 \(4\) 个白球和 \(3\) 个黑球，从中任意摸出 \(3\) 个球，其中只有一个黑球的概率是\fillinblank{}.
\end{problem}

\begin{answer}
\(\dfrac{18}{35}\)
\end{answer}

\begin{problem}
函数 \(y=\begin{cases}
2x+3,&x\leqslant0\\
x+3,&0<x\leqslant1\\
-x+5,&x>1
\end{cases}\) 的最大值是\fillinblank{}.
\end{problem}

\begin{answer}
4
\end{answer}

\begin{problem}
设 \(n\) 是一个自然数，\(\left(1+\frac{x}{n}\right)^{n} \) 的展开式中 \(x^{3} \) 的系数为 \(\frac{1}{16} \)，则 \(n=\)\fillinblank{}.
\end{problem}

\begin{answer}
4
\end{answer}

\begin{problem}
在数列 \(\{a_{n}\} \) 和 \(\{b_{n}\} \) 中， \(a_{1}=2 \)，且对任意自然数 \(n\)， \(3a_{n+1}-a_{n}=0 \)， \(b_{n} \) 是 \(a_{n} \) 与 \(a_{n+1} \) 的等差中项，则 \(\{b_{n}\} \) 的各项和是\fillinblank{}.
\end{problem}

\begin{answer}
2
\end{answer}

\begin{problem}
与椭圆 \(\frac{x^2}{49} + \frac{y^2}{24} = 1\) 有相同焦点且以 \(y = \pm \frac{4}{3} x\) 为渐近线的双曲线方程是\fillinblank{}.
\end{problem}

\begin{answer}
\(\dfrac{x^2}{9}-\dfrac{y^2}{16}=1\)
\end{answer}

\section{单选题}

\begin{problem}
下列函数中，周期为 \(\frac{\pi}{2}\) 的偶函数是（\quad）

\choices
    {\(y = \sin 4x \)}
    {\(y = \cos^2 2x - \sin^2 2x \)}
    {\(y = \tan 2x \)}
    {\(y = \cos 2x \)}
\end{problem}

\begin{answer}
B
\end{answer}

\begin{problem}
若 \(0 < a < 1\)，则函数 \(y = \log_a(x + 5)\) 的图象不经过（\quad）

\choices
    {第一象限}
    {第二象限}
    {第三象限}
    {第四象限}
\end{problem}

\begin{answer}
A
\end{answer}

\begin{problem}
在下列命题中，假命题是（\quad）

\choices
    {若平面 \(\alpha \) 内的一条直线 \(l\) 垂直于平面 \(\beta \) 内的任一直线，则 \(\alpha \perp \beta \)}
    {若平面 \(\alpha \) 内的任一直线平行于平面 \(\beta \)，则 \(\alpha \parallel \beta \)}
    {若平面 \(\alpha \perp \beta \)，任取直线 \(l \subset \alpha \)，则必有 \(l \perp \beta \)}
    {若平面 \(\alpha \parallel \beta\)，任取直线 \(l \subset \alpha\)，则必有 \(l \parallel \beta\)}
\end{problem}

\begin{answer}
C
\end{answer}

\begin{problem}
设全集为 \(\R\)，\(A=\{x\mid x^{2}-5x-6>0\}\)，\(B=\{x\mid \lvert x-5\rvert<a\}\)（\(a\) 是常数），且 \(11\in B\)，则（\quad）

\choices
    {\(\overline{A} \cup B = \R\)}
    {\(A \cup \overline{B} = \R\)}
    {\(\overline{A} \cup \overline{B} = \R\)}
    {\(A \cup B = \R\)}
\end{problem}

\begin{answer}
D
\end{answer}

\begin{problem}
设 \(a\)，\(b\)，\(c\) 分别是 \(\triangle ABC\) 中 \(\angle A\)，\(\angle B\)，\(\angle C\) 所对边的边长，则直线 \(\sin A \cdot x + ay + c = 0\) 与 \(bx - \sin B \cdot y + \sin C = 0\) 的位置关系是（\quad）

\choices
    {平行}
    {重合}
    {垂直}
    {相交但不垂直}
\end{problem}

\begin{answer}
C
\end{answer}

\section{解答题}

\begin{problem}
设 \(\alpha \) 是第二象限角，\(\sin\alpha=\frac{3}{5} \)，求 \(\sin\left(\frac{37}{6}\pi-2\alpha\right) \) 的值.
\end{problem}

\begin{solution}
因为 \(\alpha\) 是第二象限角，\(\sin\alpha=\dfrac{3}{5}\)，所以 \(\cos\alpha=-\dfrac{4}{5}\).

\(\sin2\alpha=2\sin\alpha\cos\alpha=-\dfrac{24}{25}\)，\(\cos2\alpha=1-2\sin^2\alpha=\dfrac{7}{25}\).

\[
\sin\left(\frac{37}{6}\pi-2\alpha\right)=\sin\left(\frac{\pi}{6}-2\alpha\right)=\sin\frac{\pi}{6}\cos2\alpha-\cos\frac{\pi}{6}\sin2\alpha=\frac{7+24\sqrt{3}}{50}.
\]
\end{solution}

\begin{problem}
已知复数 \(z_{1} \) 满足 \((z_{1}-2)\i=1+\i \)，复数 \(z_{2} \) 的虚部为 \(2\)，且 \(z_{1}\cdot z_{2} \) 是实数，求复数 \(z_{2} \) 的模.
\end{problem}

\begin{solution}
由 \((z_1-2)\i=1+\i\)，得 \(z_1\i=1+3\i\)，所以 \(z_1=\dfrac{1+3\i}{\i}=3-\i\).

设 \(z_2=t+2\i\)（\(t\in\R\)），则

\[
z_1z_2=(3-\i)(t+2\i)=(3t+2)+(6-t)\i.
\]

因为 \(z_1z_2\in\R\)，所以 \(6-t=0\)，得 \(t=6\).

所以 \(z_2=6+2\i\)，\(\lvert z_2\rvert=\sqrt{6^2+2^2}=2\sqrt{10}\).
\end{solution}

\begin{problem}
直四棱柱 \(ABCD - A_{1}B_{1}C_{1}D_{1} \) 的高为 \(6\)，底面是边长为 \(4\)， \(\angle DAB = 60^{\circ} \) 的菱形，\(AC\) 与 \(BD\) 相交于 \(O\)，\(A_{1}C_{1}\) 与 \(B_{1}D_{1}\) 相交于 \(O_{1}\)，\(E\) 是 \(O_{1}A\) 的中点.

\begin{enumerate}
\item 求二面角 \(O_{1}-BC-D \) 的大小（用反三角函数表示）.
\item 分别以射线 \(OA\)，\(OB\)，\(OO_{1}\) 为 \(x\) 轴，\(y\) 轴，\(z\) 轴的正半轴建立空间直角坐标系，求点 \(B_{1}\)，\(D_{1}\)，\(E\) 的坐标，并求异面直线 \(OB_{1}\) 与 \(D_{1}E\) 所成角的大小（用反三角函数表示）.
\end{enumerate}
\end{problem}

\begin{solution}
\begin{enumerate}
\item 因为 \(DD_1\perp\) 底面 \(ABCD\)，\(O_1O\parallel D_1D\)，所以 \(O_1O\perp\) 底面 \(ABCD\). 过 \(O\) 作 \(OF\perp BC\) 交 \(BC\) 于 \(F\)，连接 \(O_1F\)，根据三垂线定理得 \(O_1F\perp BC\)，所以 \(\angle O_1FO\) 是二面角 \(O_1-BC-D\) 的平面角.

\(OC=OA=4\cos30^{\circ}=2\sqrt{3}\)，\(OB=OD=2\)，\(OF=\dfrac{OC\cdot OB}{BC}=\sqrt{3}\).

\(\tan\angle O_1FO=\dfrac{6}{\sqrt{3}}=2\sqrt{3}\)，所以二面角 \(O_1-BC-D\) 的大小为 \(\arctan 2\sqrt{3}\).

\item 由条件得 \(O(0,0,0)\)，\(B_1(0,2,6)\)，\(D_1(0,-2,6)\)，\(E(\sqrt{3},0,3)\).

\(\overrightarrow{OB_1}=(0,2,6)\)，\(\overrightarrow{D_1E}=(\sqrt{3},2,-3)\).

设向量 \(\overrightarrow{OB_1}\) 与 \(\overrightarrow{D_1E}\) 的夹角为 \(\alpha\)，则

\[
\cos\alpha=\frac{\overrightarrow{OB_1}\cdot\overrightarrow{D_1E}}{\lvert\overrightarrow{OB_1}\rvert\lvert\overrightarrow{D_1E}\rvert}=-\frac{7\sqrt{10}}{40}.
\]

由此可得异面直线 \(OB_1\) 与 \(D_1E\) 所成的角为 \(\arccos\dfrac{7\sqrt{10}}{40}\).
\end{enumerate}
\end{solution}

\begin{problem}
\begin{enumerate}
\item 设动点 \(M\) 到定点 \(D(2,0) \) 的距离与到 \(y\) 轴的距离相等，求点 \(M\) 的轨迹 \(C\) 的方程.
\item 过点 \(D(2,0)\) 的直线 \(l\) 交上述轨迹 \(C\) 于 \(P\)，\(Q\) 两点，\(E\) 点坐标是 \((1,0)\)，若 \(\triangle EPQ\) 的面积为 \(4\)，求直线 \(l\) 的倾斜角 \(\alpha\) 的值.
\end{enumerate}
\end{problem}

\begin{solution}
\begin{enumerate}
\item 根据抛物线的定义，轨迹 \(C\) 是以 \(D\) 为焦点、\(y\) 轴为准线的抛物线，其顶点坐标是 \((1,0)\)，\(p=2\)，开口向右.

所以轨迹 \(C\) 的方程为 \(y^2=4(x-1)\).

设动点 \(M\) 的坐标为 \((x,y)\)，则 \(\sqrt{(x-2)^2+y^2}=\lvert x\rvert\). 两边平方得 \(x^2-4x+4+y^2=x^2\)，所以轨迹 \(C\) 的方程为 \(y^2=4(x-1)\).

\item

解法一：设直线 \(l\) 的方程为 \(y=k(x-2)\)，因 \(l\) 与抛物线有两个交点，故 \(k\ne0\). 得 \(x=\dfrac{y}{k}+2\)，代入 \(y^2=4(x-1)\) 得

\[
y^2-\frac{4}{k}y-4=0.
\]

\(\Delta=\dfrac{16}{k^2}+16>0\) 恒成立，记这个方程的两实根为 \(y_1,y_2\)，则

\[
\lvert PQ\rvert=\sqrt{1+\frac{1}{k^2}}\lvert y_1-y_2\rvert=\sqrt{1+\frac{1}{k^2}}\sqrt{(y_1+y_2)^2-4y_1y_2}=\frac{4(k^2+1)}{k^2}.
\]

又点 \(E\) 到直线 \(l\) 的距离

\[
d=\frac{\lvert k\cdot1-0-2k\rvert}{\sqrt{k^2+1}}=\frac{\lvert k\rvert}{\sqrt{k^2+1}}.
\]

所以 \(\triangle EPQ\) 的面积为

\[
S_{\triangle EPQ}=\frac{1}{2}\lvert PQ\rvert\cdot d=\frac{2\sqrt{k^2+1}}{\lvert k\rvert}.
\]

由 \(\dfrac{2\sqrt{k^2+1}}{\lvert k\rvert}=4\) 解得 \(k^2=\dfrac{1}{3}\)，所以 \(k=\pm\dfrac{\sqrt{3}}{3}\).

因此 \(\alpha=\dfrac{\pi}{6}\) 或 \(\alpha=\dfrac{5\pi}{6}\).

解法二：设直线 \(l\) 的方程为 \(y=k(x-2)\)，代入 \(y^2=4(x-1)\)，得

\[
k^2x^2-(4k^2+4)x+4k^2+4=0.
\]

因 \(l\) 与抛物线有两个交点，故 \(k\ne0\)，而 \(\Delta=16(k^2+1)>0\) 恒成立，记这个方程的两个实根为 \(x_1,x_2\). 因抛物线 \(y^2=4(x-1)\) 的焦点为 \(D(2,0)\)，准线为 \(x=0\)，所以

\[
\lvert PQ\rvert=x_1+x_2=\frac{4(k^2+1)}{k^2}.
\]

其余同解法一.
\end{enumerate}
\end{solution}

\begin{problem}
设某物体一天中的温度 \(T\) 是时间 \(t\) 的函数：\(T(t)=at^3+bt^2+ct+d\)（\(a\ne0\)），其中温度的单位是 \(^{\circ}\mathrm{C}\)，时间的单位是小时，\(t=0\) 表示 \(12{:}00\)，\(t\) 取正值表示 \(12{:}00\) 以后. 若测得该物体在 \(8{:}00\) 的温度为 \(8\mathrm{~^{\circ}C}\)，\(12{:}00\) 的温度为 \(60\mathrm{~^{\circ}C}\)，\(13{:}00\) 的温度为 \(58\mathrm{~^{\circ}C}\)，且已知该物体的温度在 \(8{:}00\) 和 \(16{:}00\) 有相同的变化率.

\begin{enumerate}
\item 写出该物体的温度 \(T\) 关于时间 \(t\) 的函数关系式.
\item 该物体在 \(10{:}00\) 到 \(14{:}00\) 这段时间中（包括 \(10{:}00\) 和 \(14{:}00\)），何时温度最高并求出最高温度.
\item 如果规定一个函数 \(f(x) \) 在 \([x_1, x_2]\)（\(x_1 < x_2\)）上函数值的平均为 \(\frac{1}{x_2 - x_1} \int_{x_1}^{x_2} f(x) dx\)，求该物体在 \(8{:}00\) 到 \(16{:}00\) 这段时间内的平均温度.
\end{enumerate}
\end{problem}

\begin{solution}
\begin{enumerate}
\item 由条件得 \(T(0)=60\)，\(T(-4)=8\)，\(T(1)=58\)，\(T'(-4)=T'(4)\)，则 \(d=60\)，\(b=0\)，\(a=1\)，\(c=-3\).

因此，温度函数 \(T(t)=t^3-3t+60\).

\item \(T'(t)=3t^2-3=3(t-1)(t+1)\).

当 \(t\in(-2,-1)\cup(1,2)\) 时，\(T'(t)>0\)；当 \(t\in(-1,1)\) 时，\(T'(t)<0\).

因此，函数 \(T(t)\) 在 \((-2,-1)\) 上单调递增，在 \((-1,1)\) 上单调递减，在 \((1,2)\) 上单调递增，即 \(t=-1\) 是极大值点.

由于 \(T(-1)=T(2)=62\)，所以在 \(10{:}00\) 到 \(14{:}00\) 这段时间中，该物体在 \(11{:}00\) 和 \(14{:}00\) 的温度最高，最高温度为 \(62\mathrm{~^{\circ}C}\).

\item 根据定义，平均温度为

\[
\frac{1}{4-(-4)}\int_{-4}^{4}T(t)\mathrm{d}t=\frac{1}{8}\int_{-4}^{4}(t^3-3t+60)\mathrm{d}t=60.
\]

即该物体在 \(8{:}00\) 到 \(16{:}00\) 这段时间内的平均温度为 \(60\mathrm{~^{\circ}C}\).
\end{enumerate}
\end{solution}

\begin{problem}
若 \(A_{n}\) 和 \(B_{n}\) 分别表示数列 \(\{a_{n}\}\) 和 \(\{b_{n}\}\) 前 \(n\) 项的和，对任意正整数 \(n\)，\(a_{n} = -\frac{2n+3}{2}\)，\(4B_{n} - 12A_{n} = 13n\).

\begin{enumerate}
\item 求数列 \(\{b_{n}\}\) 的通项公式.
\item 设有抛物线列 \(C_{1}\)，\(C_{2}\)，\(\cdots\)，\(C_{n}\)，\(\cdots\)，抛物线 \(C_{n}\)（\(n \in \N^{*}\)）的对称轴平行于 \(y\) 轴，顶点为 \((a_{n}, b_{n})\)，且通过点 \(D_{n}(0, n^{2} + 1)\)，过点 \(D_{n}\) 且与抛物线 \(C_{n}\) 相切的直线斜率为 \(k_{n}\)，求极限 \(\displaystyle\lim_{n \to \infty} \frac{k_1 + k_2 + k_3 + \cdots + k_n}{a_n b_n}\).
\item 设集合 \(X = \{x \mid x = 2a_n, n \in \N^{*}\}\)，\(Y = \{y \mid y = 4b_n, n \in \N^{*}\}\). 把 \(X \cap Y\) 中小于 \(-100\) 的元素从大到小排列构成数列 \(\{c_n\}\)，求 \(c_n\) 的通项公式.
\end{enumerate}
\end{problem}

\begin{solution}
\begin{enumerate}
\item 当 \(n\geqslant2\) 时，

\[
\begin{cases}
4B_n-12A_n=13n,\\
4B_{n-1}-12A_{n-1}=13(n-1).
\end{cases}
\]

得 \(4b_n-12a_n=13\)，\(4b_n=12a_n+13=-12n-5\)，所以 \(b_n=-\dfrac{12n+5}{4}\).

而 \(4b_1-12a_1=13\)，可得 \(b_1=-\dfrac{17}{4}\)，所以 \(\{b_n\}\) 的通项公式是 \(b_n=-\dfrac{12n+5}{4}\).

\item 抛物线 \(C_n\) 的方程为

\[
y=a\left(x+\frac{2n+3}{2}\right)^2-\frac{12n+5}{4}.
\]

因为 \(D_n(0,n^2+1)\) 在此抛物线上，所以 \(a=1\)，因此 \(C_n\) 的方程为

\[
y=\left(x+\frac{2n+3}{2}\right)^2-\frac{12n+5}{4},
\]

即

\[
y=x^2+(2n+3)x+n^2+1.
\]

所以 \(y'=2x+(2n+3)\)，\(D_n\) 处切线斜率 \(k_n=2n+3\).

\[
\lim_{n\to\infty}\frac{k_1+k_2+k_3+\cdots+k_n}{a_nb_n}=\lim_{n\to\infty}\frac{n^2+4n}{\left(-n-\frac{3}{2}\right)\left(-3n-\frac{5}{4}\right)}=\frac{1}{3}.
\]

\item 对任意 \(n\in\N^{*}\)，\(2a_n=-2n-3\)，\(4b_n=-12n-5=-2(6n+1)-3\in X\)，所以 \(Y\subseteq X\)，故 \(X\cap Y=Y\).

\(4b_n=-12n-5<-100\)，故 \(n>7\dfrac{11}{12}\)，\(n\in\N^{*}\)，满足条件的最小自然数 \(n=8\).

所以 \(c_1=4b_8=-101\)，\(c_n=-101-12(n-1)=-12n-89\).
\end{enumerate}
\end{solution}
